\documentclass[11pt]{article}
\usepackage[margin=1in]{geometry}
\usepackage{amsmath,amssymb,amsthm,booktabs,array}
\usepackage[T1]{fontenc}
\usepackage[colorlinks=true,urlcolor=blue,citecolor=blue,linkcolor=blue]{hyperref}
\newtheorem{theorem}{Theorem}
\newtheorem{proposition}[theorem]{Proposition}
\newtheorem{corollary}[theorem]{Corollary}
\title{Integral-product extensions and an alternative proof\protect\\
of Dong's (2026) theorem on\protect\\
OEIS A389472, revision 7, COMMENT 3}
\author{Haobo Ma\protect\\\small ChronoAI Pte Ltd, Singapore
\and Wenlin Zhang\protect\\\small National University of Singapore\protect\\\small ChronoAI Pte Ltd, Singapore}
\date{November 10, 2026\protect\\\small Licensed under CC BY 4.0}
\begin{document}
\maketitle
\begin{abstract}
For Paul D. Hanna's OEIS A389472, revision 7 (2025), COMMENT 3,
proved by Horace Dong (2026), we give an alternative proof of
$a(3n-1)\equiv0\pmod3$ for every $n>1$, and integral-product extensions.
The recurrence proof combines coefficient extraction with the
least base-three digits of its binomial coefficients.
An infinite-product identity relates this sequence to Hanna's A350432;
it also proves the companion parity conjecture and the stated radius of
convergence. A single product calculation also establishes the functional,
inverse and radius identities of ten cross-referenced sequences, and identifies
a missing series reversion in A273162's product formula.
\end{abstract}

\section{The normalized series and the conjecture}
Hanna's entry A389472~\cite{hanna472}, revision 7, lists two congruence
conjectures. Its COMMENT 3 states
``Conjecture: a(3*n-1) = 0 (mod 3) for n > 1.''
Theorem~\ref{thm:main} proves this statement for the series
\[
 A(x)=\sum_{m\ge1}a_mx^m
     =x+x^2+2x^3+3x^4+6x^5+\cdots,
 \qquad q(x)=x^2+x^3,
\]
whose defining equation is
\begin{equation}\label{eq:def}
 A(q(x))=x^2(1+A(x)).
\end{equation}
The third COMMENT is the modulo-three conjecture proved below;
the second is its parity companion. The initial condition $a_2=1$
is part of the definition: equation~\eqref{eq:def} alone does not fix it.
We put $a_0=0$ throughout.

Dong~\cite{dong}, Section 3, Theorem 4, proves the modulo-three assertion
for every integral solution of~\eqref{eq:def}, using formal derivatives
in characteristic three. Perez Fontelles~\cite{perez}, Section 4,
proves the parity assertion; his Section 2 also gives the normalized
recurrence and uniqueness. The proof below is an alternative
binomial argument. The relation to Hanna's integral product supplies
the classification and the further product consequences.

For $m\ge3$, coefficient extraction at degree $m+2$ gives
\begin{equation}\label{eq:rec}
 a_m=\sum_{k=0}^{\lfloor(m+2)/2\rfloor}
       a_k\binom{k}{m+2-2k},
\end{equation}
where $\binom{k}{j}=0$ for $j>k$. Indeed,
$q(x)^k=x^{2k}(1+x)^k$, so its coefficient at degree $m+2$
is the displayed binomial coefficient. Every index on the right is
smaller than $m$ for $m\ge3$. Thus $a_1=a_2=1$ and
\eqref{eq:rec} determine the series uniquely over the integers.
The coefficients at degrees $0,1,2,3,4$ of \eqref{eq:def}
are consistent with these initial values, so the recurrence also establishes
existence of the normalized formal series.

\begin{theorem}\label{thm:main}
For the normalized series \eqref{eq:def},
\[
 3\mid a_{3n-1}\qquad\text{for every integer }n>1.
\]
\end{theorem}
\begin{proof}
We first record the elementary digit calculation needed in the induction.
Write $k=3u+r$ and $j=3v+s$, where $0\le r,s<3$. In
$\mathbb F_3[t]$,
\[
 (1+t)^k=(1+t^3)^u(1+t)^r,
 \qquad
 \binom{k}{j}\equiv\binom{u}{v}\binom{r}{s}\pmod3.
\]
The coefficient formula follows by comparing the powers whose exponent is
$3v+s$: the second factor has degree at most two. In particular,
$s>r$ implies $3\mid\binom{k}{j}$.

We prove $3\mid a_m$ for all $m\ge5$ with $m\equiv2\pmod3$
by strong induction on $m$. Set $d=m+2\equiv1\pmod3$ and
apply~\eqref{eq:rec}. In a summand let $j=d-2k$.
If $k\equiv0\pmod3$, then $j\equiv1\pmod3$;
if $k\equiv1\pmod3$, then $j\equiv2\pmod3$.
In both cases the digit calculation makes the binomial factor divisible
by three. If $k\equiv2\pmod3$ and $k\ge5$, the induction
hypothesis makes $a_k$ divisible by three, since $k<m$.
The remaining possibility is $k=2$, but then
$j=m-2\ge3>2$, so the binomial coefficient is zero.
Every summand is therefore divisible by three. This includes the initial
case $m=5$, without a further base assumption. Taking $m=3n-1$ proves
the assertion. The restriction $n>1$ is necessary because $a_2=1$.
\end{proof}

\section{A product identity and the companion statements}
Let $q_0(x)=x$ and $q_{j+1}(x)=q(q_j(x))$. Define
\begin{equation}\label{eq:product}
 P(x)=\prod_{j\ge0}(1+q_j(x)),\qquad H(x)=xP(x).
\end{equation}
Here $\operatorname{ord}$ denotes the least exponent with a nonzero coefficient.
These products are well defined as formal power series because
$\operatorname{ord}(q_j)=2^j$. Hanna's earlier entry
A350432~\cite{hanna432}, FORMULA (3), gives exactly $H$:
its factor $F(j+1)$ is $1+q_j$. Thus the product used here is
Hanna's construction, not a new product representation of A350432.

\begin{proposition}\label{prop:product}
The normalized series satisfies
\begin{equation}\label{eq:bridge}
 A(x)=x+H(x)^2=x+x^2P(x)^2.
\end{equation}
Consequently $2\mid a_{2n-1}$ for every $n>1$, and
$6\mid a_{6n-1}$ for every $n\ge1$.
Its radius of convergence is $r=(\sqrt5-1)/2$.
\end{proposition}
\begin{proof}
Shifting the factors in~\eqref{eq:product} gives
$P(x)=(1+x)P(q(x))$, hence
\[
 H(q(x))=q(x)P(q(x))=xH(x).
\]
Therefore the series $B(x)=x+H(x)^2$ has coefficients
$[x]B=[x^2]B=1$ and satisfies
\[
 B(q(x))=q(x)+x^2H(x)^2=x^2(1+B(x)).
\]
Uniqueness proves~\eqref{eq:bridge}. Over $\mathbb F_2$ the square
of any power series has zero coefficients in odd degrees. Thus all odd
coefficients of $A-x$ are even, proving the second COMMENT of A389472.
An index $6n-1$ is both odd and congruent to two modulo three, and is at
least five. Theorem~\ref{thm:main} and parity give the factor six.

For completeness, the same product proves the first COMMENT, the stated
radius. Fix $0<\rho<r$ and put $c=\rho+\rho^2<1$.
For $|z|\le\rho$, positivity of the coefficients of $q_j$ gives
$|q_j(z)|\le q_j(\rho)$, and
\[
 0\le q_{j+1}(\rho)
   =q_j(\rho)^2(1+q_j(\rho))
   \le c\,q_j(\rho)\le\rho c^{j+1}.
\]
The sum of these bounds converges, so the product converges uniformly
on $|z|\le\rho$ and defines an analytic function there.
Thus $A$ has radius at least $r$. On the other hand $q_j(r)=r$ for
every $j$. Each finite product has nonnegative coefficients and is
coefficientwise bounded by $P$; by~\eqref{eq:bridge},
\[
 \sum_{m\ge1}a_mr^m
 \ \ge\ r+r^2(1+r)^{2(J+1)}\qquad(J\ge0).
\]
The right side is unbounded. The series diverges at $r$, proving that
its radius is exactly $r$.
\end{proof}

\begin{corollary}\label{cor:normalization}
Every integer-coefficient solution of~\eqref{eq:def} with constant
coefficient zero is $A_b(x)=x+bH(x)^2$ for a unique integer $b$.
Both congruences and their factor-six consequence hold for every $b$.
For $b\ne0$ the radius is $r$; for $b=0$ the solution is the polynomial $x$.
\end{corollary}
\begin{proof}
The coefficient at degree two in~\eqref{eq:def} forces $[x]A_b=1$;
the degree-four equation leaves $[x^2]A_b=b$ free. Substitution proves
that $x+bH^2$ is a solution, and~\eqref{eq:rec} proves uniqueness.
For $m\ge2$ its coefficient is $b a_m$, so both divisibilities follow.
For $b\ne0$, multiplying the tail of $A$ by $b$ does not change its radius.
\end{proof}

The same bounds give radius $r$ for $H$: its finite products at $r$
are $r(1+r)^{J+1}$ and are unbounded. Moreover $H'(0)=1$, so its
formal compositional inverse $S$ exists, and $H(q(x))=xH(x)$ gives
$S(xH(x))=q(x)$. This is FORMULA (2) of A350432; the equation for
$H$ above is its FORMULA (1).

The remaining two numbered formulas in A389472 follow directly from
\eqref{eq:def}. Since $A'(0)=1$, its formal compositional inverse
$R$ exists. Substituting $x=R(t)$ gives FORMULA (2):
\[
 A\bigl((1+R(t))R(t)^2\bigr)=(1+t)R(t)^2.
\]
Applying~\eqref{eq:def} twice gives
\[
 A(q(q(x)))=q(x)^2\bigl(1+x^2+x^2A(x)\bigr).
\]
Now $q(q(x))=x^4+2x^5+2x^6+3x^7+3x^8+x^9$ and
$x^2q(x)^2=(x^3+x^4)^2$. Rearranging in the Laurent-series
field gives precisely FORMULA (3):
\[
 A(x)=\frac{A(x^4+2x^5+2x^6+3x^7+3x^8+x^9)}{(x^3+x^4)^2}
       -\frac{1+x^2}{x^2}.
\]

\section{Cross-referenced product families}
The same factor-shift calculation proves the following general statement;
it is not posed as a conjecture in the cited entries. All products and
inverses in this section are formal series in $\mathbb Z[[x]]$.

\begin{proposition}\label{prop:family}
Let $d\ge2$ and $s\ge1$ be integers, let $\varepsilon\in\{1,-1\}$,
and put
\[
 \begin{gathered}
 Q(x)=x^d(1+\varepsilon x)^s,\quad Q_0(x)=x,\quad Q_{j+1}(x)=Q(Q_j(x)),\\
 P_Q(x)=\prod_{j\ge0}(1+\varepsilon Q_j(x)).
 \end{gathered}
\]
For $t\in\{1,s\}$, the series $U_t(x)=xP_Q(x)^t$ is the unique
series with constant coefficient zero and linear coefficient one satisfying
\begin{equation}\label{eq:family}
 U_t(Q(x))=x^{d-1}(1+\varepsilon x)^{s-t}U_t(x).
\end{equation}
Its compositional inverse $T_t$ satisfies
\[
 T_t\bigl(x^{d-1}(1+\varepsilon x)^{s-t}U_t(x)\bigr)=Q(x).
\]
When $\varepsilon=1$, both $U_1$ and $U_s$ have radius $r_{d,s}$,
the unique positive solution of $r_{d,s}^{d-1}(1+r_{d,s})^s=1$.
\end{proposition}
\begin{proof}
The order of $Q_j$ is $d^j$, so each coefficient of the product
stabilizes. The identity $P_Q(x)=(1+\varepsilon x)P_Q(Q(x))$
gives~\eqref{eq:family}. If two normalized solutions differ first at
degree $m\ge2$, their difference after substitution has order $dm$,
but the right side has order $m+d-1$. These are unequal, proving
uniqueness. Since $U_t'(0)=1$, applying $T_t$ proves the inverse identity.

For $\varepsilon=1$ and $0<\rho<r_{d,s}$, set
$c=\rho^{d-1}(1+\rho)^s<1$. Positivity gives
$|Q_j(z)|\le Q_j(\rho)\le\rho c^j$ on $|z|\le\rho$.
Thus the product converges uniformly there. At $r_{d,s}$ every
iterate is $r_{d,s}$ and its partial products for $U_t$ equal
$r_{d,s}(1+r_{d,s})^{t(J+1)}$. Nonnegative coefficients imply
divergence at that point, proving the asserted radius.
\end{proof}

\begin{corollary}\label{cor:positive}
For $s=2,3,4,5,6$, put $Q=x^2(1+x)^s$, $U=xP_Q$ and $V=xP_Q^s$.
Then $U$ is respectively the generating function of A369545--A369549,
and $V$ that of A369552--A369556~\cite{hanna-family}. Their numbered
functional, inverse and product identities hold, and
\[
 U(x)^s=x^{s-1}V(x),\qquad
 \operatorname{rad}(U)=\operatorname{rad}(V)=r_{2,s}.
\]
\end{corollary}
\begin{proof}
Proposition~\ref{prop:family} gives $U(Q)=x(1+x)^{s-1}U$,
$V(Q)=xV$ and the inverse identities. Its factors $F(j+1)=1+Q_j$
satisfy $F(j+2)=1+(F(j+1)-1)^2F(j+1)^s$, exactly the cited
product recurrence. Uniqueness identifies the series; the power identity
is immediate from their products, and the proposition supplies the radii.
\end{proof}

\begin{corollary}\label{cor:signed}
For $Q=x^d(1-x)$ put $W_d=x\prod_{j\ge0}(1-Q_j)$ and let $Z_d$
be its compositional inverse. Then
\[
 W_d(Q)=x^{d-1}W_d,\qquad Z_d(x^{d-1}W_d(x))=x^d-x^{d+1}.
\]
The series of A273162 is $Z_2$, that of A350433 is $W_3$, and
that of A350434 is $Z_4$~\cite{hanna-signed}.
In A273162, revision 12, the second FORMULA's product is $W_2$,
the inverse series, rather than its stated generating function $Z_2$.
The corrected formula is $Z_2=W_2^{\langle-1\rangle}$.
\end{corollary}
\begin{proof}
Use $(d,s,\varepsilon,t)=(d,1,-1,1)$ in
Proposition~\ref{prop:family}. The factors now satisfy
$F(j+2)=1-(1-F(j+1))^dF(j+1)$, as in the three entries.
For A273162 the equations $Z_2(xW_2)=Q$ and $Z_2(W_2)=x$
also give, by substituting $x=Z_2(y)$,
$Z_2(yZ_2(y))=Z_2(y)^2-Z_2(y)^3$, its first FORMULA.
The second FORMULA literally asserts $Z_2=x\prod_{j\ge0}(1-Q_j)$.
Its two sides have second coefficients $+1$ and $-1$ respectively,
so that assertion fails. Its EXAMPLE field correctly assigns the product
to the inverse $W_2$, consistently with the replacement above.
The same inverse identity identifies $W_3$ and $Z_4$ with the
normalized series in their entries.
\end{proof}

Taking $(d,s,\varepsilon)=(2,1,1)$ also proves the defining
identity and reversion formula of A350431~\cite{hanna431}: its
series is $H^{\langle-1\rangle}$, so
$H^{\langle-1\rangle}(xH(x))=x^2+x^3$.
The positive family extends the radius argument; no analytic-radius
assertion for the signed family is needed here.

\section{Coverage and verification}
The following table identifies the mathematical statements in the posing
entry. The numerical values in its EXAMPLE field are illustrations rather
than additional claims proved in this note.
\begin{center}
\small
\begin{tabular}{@{}p{0.43\linewidth}p{0.51\linewidth}@{}}
\toprule
A389472, revision 7 & Coverage \\
\midrule
COMMENT 3: modulo-three conjecture & Reproved in Theorem~\ref{thm:main};
prior proof: Dong (September 26, 2026), Theorem 4. \\
COMMENT 2: parity conjecture & Reproved in Proposition~\ref{prop:product};
prior proof: Perez Fontelles (August 31, 2026), Section 4. \\
COMMENT 1: radius $r$ & Proved in Proposition~\ref{prop:product}. \\
FORMULA (1): defining equation & Existence and uniqueness with $a_2=1$;
prior argument: Perez Fontelles, Section 2. \\
FORMULA (2): inverse-series identity & Proved by substitution above. \\
FORMULA (3): iterated substitution & Proved by iteration above. \\
Cross-reference A350432 & Hanna's product (3) is cited; formulas (1), (2)
and its radius follow above. The relation to $A$ is~\eqref{eq:bridge}. \\
General second coefficient $b$ & Not posed; proved in
Corollary~\ref{cor:normalization}. \\
\bottomrule
\end{tabular}
\end{center}
The cross-references of A350432 have the following additional coverage.
These are proofs of Hanna's stated identities, not claims that the identities
originate here. Entries with no numbered formulas are located by field.
\begin{center}
\small
\begin{tabular}{@{}p{0.43\linewidth}p{0.51\linewidth}@{}}
\toprule
Source item & Coverage \\
\midrule
A350431, NAME and FORMULA & Proved by reversion above. \\
A369545--A369549; A369552--A369556, COMMENT and numbered FORMULAS
& Proved in Corollary~\ref{cor:positive}; Hanna's product formulas cited.
Decimal approximations are not included. \\
A273162, NAME and FORMULA 1 & Proved in Corollary~\ref{cor:signed}. \\
A273162, revision 12, FORMULA 2 & Refuted-and-corrected:
the product is the inverse series, not the sequence's series. \\
A350433, NAME and FORMULA; A350434, NAME and FORMULA
& Proved in Corollary~\ref{cor:signed}. \\
Further decimal evaluations and extrema in EXAMPLE fields
& Not covered; these are distinct from the formal identities. \\
General $d,s,\varepsilon$ family & Not posed; proved in
Proposition~\ref{prop:family}. \\
\bottomrule
\end{tabular}
\end{center}

The pinned Lean declaration \texttt{result} in the
\href{https://github.com/the-omega-institute/trureturing/blob/074df09ddbbef1843ddb383577bfaa827f753e36/D5/S1/Recurrence/Parity/SquareCubeSubstitutionThreeIndexModThree.lean}{modulo-three source}
has the statement $\forall n\in\mathbb N,\ 1<n\Rightarrow3\mid a(3n-1)$.
The companion
\href{https://github.com/the-omega-institute/trureturing/blob/074df09ddbbef1843ddb383577bfaa827f753e36/D5/S1/Recurrence/Parity/SquareCubeSubstitutionOddIndexParity.lean}{parity source}
contains \texttt{generating\_equation}, \texttt{generating\_unique},
\texttt{coeff\_recurrence} and \texttt{hanna\_conjecture}; these identify the
normalized sequence and the parity statement. The product identity, analytic
radius, inverse-series and iteration arguments are written proofs above;
the factor-six conclusion is their elementary arithmetic corollary.

The accompanying \texttt{verify.py} compares the binomial recurrence with
the iterated product, checks the displayed source coefficients and the
congruences, and checks the product-family identities and their inverses
over explicitly bounded ranges recorded in \texttt{checks.json}.
These are finite integer computations; the all-index and analytic conclusions
have the written proofs above. The additional corollaries are not part of
the two linked Lean declarations' scope.

Dong's earlier
\href{https://github.com/Horace-Maxwell/ai4math-results/blob/4c9a4c2b6945ed4f12ba8bdd66c398869081b4f6/lean/Research/Oeis389472Mod3.lean}{\texttt{Oeis389472Mod3.integer\_conjecture}}
states the modulo-three conclusion for any integer power series satisfying
\eqref{eq:def}; unlike the normalized recurrence declaration, it does not
assume a specified second coefficient. The two scopes agree on A389472.

\begin{thebibliography}{9}
\bibitem{dong}
H. Dong, \emph{Proofs of eleven congruence conjectures of Hanna from the
OEIS, and the resolution of three more}, September 2026,
Section 3, Theorem 4; version 1.1.0, archived September 26, 2026.
\href{https://doi.org/10.5281/zenodo.22970816}{Versioned archive},
\texttt{papers/oeis-hanna/note.pdf}.
\bibitem{perez}
A. Perez Fontelles, \emph{Every odd-indexed term of OEIS A389472 after
the first is even}, August 31, 2026, Sections 2 and 4.
\href{https://github.com/astrafala/Conjectures/blob/58cf4596af4541652fb191e177ef9b7771388d53/paper-sources/01001-01500/01042-PROOF.tex}{Pinned manuscript}.
\bibitem{hanna472}
P. D. Hanna, \emph{A389472: G.f. satisfies
$A(x)=A(x^2+x^3)/x^2-1$}, The On-Line Encyclopedia of Integer Sequences,
October 15, 2025; revision 7, October 17, 2025.
\href{https://oeis.org/A389472}{Entry} and
\href{https://oeis.org/history/view?seq=A389472\&v=7}{revision 7}.
\bibitem{hanna432}
P. D. Hanna, \emph{A350432: G.f. $A(x)$ satisfies
$A(x)=A(x^2+x^3)/x$}, The On-Line Encyclopedia of Integer Sequences,
December 30, 2021; revision 13, January 26, 2024, FORMULA (3).
\href{https://oeis.org/A350432}{Entry} and
\href{https://oeis.org/history/view?seq=A350432\&v=13}{revision 13}.
\bibitem{hanna431}
P. D. Hanna, \emph{A350431}, The On-Line Encyclopedia of Integer Sequences,
December 30, 2021; revision 5.
\href{https://oeis.org/A350431}{Entry}.
\bibitem{hanna-family}
P. D. Hanna, \emph{A369545--A369549 and A369552--A369556},
The On-Line Encyclopedia of Integer Sequences, January 25--26, 2024.
\href{https://oeis.org/A369545}{A369545},
\href{https://oeis.org/A369546}{A369546},
\href{https://oeis.org/A369547}{A369547},
\href{https://oeis.org/A369548}{A369548},
\href{https://oeis.org/A369549}{A369549};
\href{https://oeis.org/A369552}{A369552},
\href{https://oeis.org/A369553}{A369553},
\href{https://oeis.org/A369554}{A369554},
\href{https://oeis.org/A369555}{A369555},
\href{https://oeis.org/A369556}{A369556}.
\bibitem{hanna-signed}
P. D. Hanna, \emph{A273162}, May 16, 2016 (revision 12,
December 30, 2021); \emph{A350433 and A350434}, December 30, 2021,
The On-Line Encyclopedia of Integer Sequences.
\href{https://oeis.org/A273162}{A273162},
\href{https://oeis.org/A350433}{A350433},
\href{https://oeis.org/A350434}{A350434}.
\end{thebibliography}
\end{document}
